\documentclass{article}
\usepackage{amsmath}
\usepackage{longtable}
\usepackage{booktabs}
\begin{document}

\section{Product Description}

The \textit{Aurora} multispectral imager provides measurements in
three spectral bands. Table~\ref{tab:bands} lists the characteristics
of each channel.

\begin{longtable}{|l|c|c|}
\caption{Spectral bands}\label{tab:bands}\\
\hline
Band & Wavelength & Resolution\\
\hline
\endfirsthead
\hline
Band & Wavelength & Resolution\\
\hline
\endhead
blue   & 443 nm & 300 m\\
green  & 555 nm & 300 m\\
red    & 670 nm & 300 m\\
\hline
\end{longtable}

\subsection{Geometric accuracy}

The geolocation error is computed as
\begin{equation}
  \epsilon = \sqrt{(\Delta x)^2 + (\Delta y)^2}
\end{equation}
where $\Delta x$ and $\Delta y$ are the along-track and cross-track
residuals, respectively.

\begin{itemize}
\item boresight calibration performed monthly
\item thermal distortion model v3
\end{itemize}

\begin{verbatim}
calibration_report_v2.txt
  band 1: offset = 0.42
  band 2: offset = 0.17
\end{verbatim}

\subsubsection{Data format}\label{sec:format}

Products are distributed in a table-driven format; see
Table~\ref{tab:fields} for the variable list.

\begin{longtable}{|l|l|}
\hline
Variable & Type\\
\hline
time       & double\\
latitude   & float32\\
longitude  & float32\\
quality    & uint8\\
\hline
\end{longtable}

\end{document}
